
\documentclass{article}
\usepackage{colortbl}
\usepackage{makecell}
\usepackage{multirow}
\usepackage{supertabular}

\begin{document}

\newcounter{utterance}

\centering \large Interaction Transcript for game `cladder', experiment `full\_v1.5\_default', episode 8669 with selfrefine{-}llama{-}3.1{-}8b{-}instruct.
\vspace{24pt}

{ \footnotesize  \setcounter{utterance}{1}
\setlength{\tabcolsep}{0pt}
\begin{supertabular}{c@{$\;$}|p{.15\linewidth}@{}p{.15\linewidth}p{.15\linewidth}p{.15\linewidth}p{.15\linewidth}p{.15\linewidth}}
    \# & \multicolumn{2}{c}{Player} && \multicolumn{2}{c}{Game Master} \\
    \hline

    \theutterance \stepcounter{utterance}  
    & & & \multicolumn{4}{p{0.6\linewidth}}{
        \cellcolor[rgb]{0.9,0.9,0.9}{
            \makecell[{{p{\linewidth}}}]{
                \texttt{\tiny{[P1$\langle$GM]}}
                \texttt{You are an expert in causal inference. The following question is not a typical commonsense query, but rather a meticulously designed question created by a professor specializing in causal inference, intended to assess the students' mastery of the course content.} \\
\texttt{Imagine a self{-}contained, hypothetical world with only the following conditions, and without any unmentioned factors or causal relationships: Jyka has a direct effect on yupt and kwox. Qwiu has a direct effect on yupt and kwox. Yupt has a direct effect on kwox. Qwiu is unobserved. The overall probability of jyka is 0\%. The probability of not jyka and kwox is 11\%. The probability of jyka and kwox is 0\%. Is the chance of kwox smaller when observing jyka?} \\
\texttt{Start your answer with 'Yes' or 'No', followed by additional reasoning or evidence to support your explanation.} \\
            }
        }
    }
     \\ \\

    \theutterance \stepcounter{utterance}  
    & & \multicolumn{4}{p{0.6\linewidth}}{
        \cellcolor[rgb]{0.95,0.95,0.95}{
            \makecell[{{p{\linewidth}}}]{
                \texttt{\tiny{[P1$\leftrightarrow$P]}}
                \texttt{{[}DRAFT{]}} \\
\texttt{No.} \\
\\ 
\texttt{To determine whether the chance of kwox is smaller when observing jyka, we need to consider the causal relationships and probabilities given in the problem.} \\
\\ 
\texttt{First, we know that Jyka has a direct effect on yupt and kwox, and Qwiu has a direct effect on yupt and kwox. This means that both Jyka and Qwiu are potential causes of yupt and kwox.} \\
\\ 
\texttt{However, the probability of Jyka is 0\%, which implies that Jyka is not a possible cause of kwox in this scenario. This is because the probability of Jyka is 0\%, making it impossible for Jyka to occur.} \\
\\ 
\texttt{The probability of not Jyka and kwox is 11\%, and the probability of Jyka and kwox is 0\%. This suggests that the occurrence of kwox is not directly related to Jyka, as the probability of both Jyka and kwox occurring is 0\%.} \\
\\ 
\texttt{Given the direct effect of yupt on kwox, we can infer that the probability of kwox is influenced by yupt. However, since Jyka has a direct effect on yupt, we can also infer that Jyka indirectly affects kwox through yupt.} \\
\\ 
\texttt{However, the key point here is that the probability of Jyka is 0\%, which means that Jyka cannot be a cause of kwox. Therefore, observing Jyka does not provide any information about the probability of kwox, as Jyka is not a possible cause of kwox.} \\
\\ 
\texttt{In conclusion, the chance of kwox is not smaller when observing Jyka, as Jyka is not a possible cause of kwox.} \\
            }
        }
    }
    & \\ \\

    \theutterance \stepcounter{utterance}  
    & & \multicolumn{4}{p{0.6\linewidth}}{
        \cellcolor[rgb]{0.95,0.95,0.95}{
            \makecell[{{p{\linewidth}}}]{
                \texttt{\tiny{[P1$\leftrightarrow$P]}}
                \texttt{{[}SELF{-}FEEDBACK{]}} \\
\texttt{FORMAT: ok} \\
\texttt{RULES: ok} \\
\texttt{STRATEGY: problem} \\
\texttt{REASON: The answer does not consider the indirect effect of Jyka on kwox through yupt, which could potentially change the probability of kwox when observing Jyka.} \\
            }
        }
    }
    & \\ \\

    \theutterance \stepcounter{utterance}  
    & \multicolumn{4}{p{0.6\linewidth}}{
        \cellcolor[rgb]{0.561,0.737,0.561}{
            \makecell[{{p{\linewidth}}}]{
                \texttt{\tiny{[P1$\rangle$GM]}}
                \texttt{No.} \\
\\ 
\texttt{The probability of Jyka is 0\%, which implies that Jyka is not a possible cause of kwox in this scenario. However, Jyka has a direct effect on yupt, and yupt has a direct effect on kwox. This means that Jyka indirectly affects kwox through yupt.} \\
\\ 
\texttt{Given the indirect effect of Jyka on kwox through yupt, we can infer that observing Jyka would provide information about the probability of yupt, which in turn affects the probability of kwox. However, the probability of not Jyka and kwox is 11\%, and the probability of Jyka and kwox is 0\%. This suggests that the occurrence of kwox is not directly related to Jyka, as the probability of both Jyka and kwox occurring is 0\%.} \\
\\ 
\texttt{Since the probability of Jyka and kwox is 0\%, observing Jyka does not provide any information about the probability of kwox, as Jyka is not a possible cause of kwox. Therefore, the chance of kwox is not smaller when observing Jyka.} \\
            }
        }
    }
    & & \\ \\

    \theutterance \stepcounter{utterance}  
    & & & \multicolumn{2}{p{0.3\linewidth}}{
        \cellcolor[rgb]{0.9,0.9,0.9}{
            \makecell[{{p{\linewidth}}}]{
                \texttt{\tiny{[GM$|$GM]}}
                \texttt{No} \\
            }
        }
    }
    & & \\ \\

    \theutterance \stepcounter{utterance}  
    & & & \multicolumn{2}{p{0.3\linewidth}}{
        \cellcolor[rgb]{0.9,0.9,0.9}{
            \makecell[{{p{\linewidth}}}]{
                \texttt{\tiny{[GM$|$GM]}}
                \texttt{no} \\
            }
        }
    }
    & & \\ \\

    \theutterance \stepcounter{utterance}  
    & & & \multicolumn{2}{p{0.3\linewidth}}{
        \cellcolor[rgb]{0.9,0.9,0.9}{
            \makecell[{{p{\linewidth}}}]{
                \texttt{\tiny{[GM$|$GM]}}
                \texttt{game\_result = WIN} \\
            }
        }
    }
    & & \\ \\

\end{supertabular}
}

\end{document}
